\section{Outline}

% Aims and research question
% % Aims are broad example improve quality of life of ALS patients.

% List of specific objectives
% -

% -

\subsection{Motivation}

Intraosseous transcutaneous amputation prostheses (ITAP) provide an alternative means of attaching artificial limbs for amputees, bypassing the role of the residual limb (stump) soft tissues in the attachment.
In the example of a leg prosthesis, we observe uneven loads, friction against socket and other abuses which lead to severe skin problems \citep{levy_skin_1980}.
This highlights one of the motivations for ITAP, shifting the load from soft tissue to the skeleton, and alleviating some of the highlighted issues.

An amputation is not just the loss of a limb, but also the loss of information one gathers with that limb.
Stump protheses, including modern robotic ones, only provide visual feedback to the user; and, in fact, here lies the motivation for this study:
One of the best tools to regain sensation is the prosthetic itself with the use of vibrotacticle feedback.
% Indeed, this leads to the possibility of using ITAP for the development of cognitive control over artificial limbs.
However, little is understood about the consequence of vibrations on either the skin-implant or the bone-implant interfaces.

Following amputation and depending on the bone quality of the residuum, ITAP can be press-fit or cemented into the remaining bone in the amputated limb – ultimately having an effect on the way vibrations are passed through ITAP and transmitted to the bone and soft tissues.
This project will investigate \textbf{the effect of press-fitting or cementing on the magnitude and nature of the forces cells encounter at the ITAP-bone and ITAP-soft tissue interfaces}.
The results of this study will feed into a larger research project to assess the effect of these forces on the biological implant interfaces between ITAP and their recipients to ensure that haptics are a safe and reliable means of controlling artificial limbs for amputees.

\subsection{Approach}
\paragraph{In silico modelling -- Finite Element analysis}
\paragraph{In vitro(?) modelling}
\begin{enumerate}
    \item 
    \item Emulate the interface between ITAP and soft tissue
\end{enumerate}

\subsection{Timing}
Rough timetable for the project
\begin{table}[h]
\centering
\begin{tabular}{clcc}
      & Task & Start & Finish \\ \hline
    \textbf{1} & \textbf{Outline} & \textbf{1 Nov} & \textbf{5 Nov} \\
    1.1 & Risk Assessment & 1 Nov & 4 Nov \\
    1.2 & Attachment research &1 Nov &4 Nov \\
    1.3 & Outline submission &5 Nov &5 Nov \\

    \textbf{2} & \textbf{Literature Review} & \textbf{8 Nov} &\textbf{ 19 Nov }\\
    2.1 & ITAP &8 Nov &12 Nov \\
    2.2 & Motors and sensors & 15 Nov & 19 Nov \\

    \textbf{3} & \textbf{FEA modelling} & \textbf{8 Nov} & \textbf{31 Dec} \\
    3.1 & Basic rod design & 8 Nov & 26 Nov \\
    3.2 & Increase complexity & 26 Nov & 31 Dec \\
    
    \textbf{4} & \textbf{Prototyping} & \textbf{31 Dec} & \textbf{18 Feb }\\
    4.1 & Measuring devices & 31 Dec & 28 Jan \\
    4.2 & Vibration motors & 31 Dec & 28 Jan \\
    4.3 & Emulation of skin-ITAP interface & 31 Dec & 18 Feb \\
    4.4 & Prototype & 18 Feb & 18 Feb \\
    
    \textbf{5} & \textbf{Project Talks} & \textbf{14 Mar} & \textbf{1 Apr} \\
    5.1 & Preparation & 14 Mar & 31 Mar\\
    5.2 & Project talks day & 1 Apr & 1 Apr\\

    \textbf{6} & \textbf{Final report} & \textbf{31 Dec} & \textbf{3 May} \\
    6.1 & Writing & 31 Dec & 10 Mar \\
    6.2 & Revision by supervisors & 10 Mar & 15 Apr \\
    6.3 & Final report due & 3 May & 3 May \\
\end{tabular}
\end{table}